%\newpage

%       Perhaps something more about the historical nature of the work. 
%       I need to not only address \textit{what} happens, but also why. Under what conditions do LMs exhibit human like behavior?
%       No YES \ NO research questions.
%       Perhaps occilate between terms and use synonoms, "cooperative interdependence"
%       Think about / through my thoughts with the Peter Principal. What are the limit conditions? Perhaps emphasize more of the connections from the popular concepts to the newer stuff.
%       Going one level out to say something about what the LMs actually value? Is it the PhD OR is it the skakeboarding?
%       In the literature gap, talk about how it's normative. Why now? Why hasn't AI systems been 
%       What art work is like this? Trevor Packland image before 2020? Be explicit that this isn't the first time that it has been addressed. ImageNet Roulet. 
%       Should I narrow down my example of hierarchy to several examples. 

% The playground offers a perfect introduction to the nearly universal experience to the dynamics of social organization. As \citet[p. 1]{Ridgeway-2019-Status} observes, ``We see status virtually everywhere in social life, \textit{if we think to look for it} [emphasis added].'' At first, as children stepping onto the playground, we may not immediately recognize the ubiquity of status. Yet with closer attention, we begin to notice how it infuses everyday interactions, from the kind of car we were dropped off in that morning, to the clothes we wear in class, and to the people we associate with on the blacktop. Our later life is likewise scored by status, from the college we attend, the neighborhood we live in, or our occupation. The organizations we work for themselves carry varying degrees of status, something we become acutely aware of when describing our affiliations to others. Consider the social implications of answering the casual question at a gathering, ``So, what do you do?'' Beyond occupational ties, status is also bound up with our social identities, such as gender, race, ethnicity, age, and class background. Moreover, in nearly every group we enter, implicit or explicit hierarchies quickly emerge, positioning us higher or lower within the social order.

\section{Literature Review}

\subsection{What is Status?}

Status, defined simply, is the position in a social hierarchy that results from accumulated acts of social esteem, honor, and respect \citep{Goode-1978-Celebration, Whyte-1943-Street, Ridgeway-1987-Nonverbal, Ridgeway-2019-Status}. \citet[p. 2]{Ridgeway-2019-Status} argues that ``status evaluations and status hierarchies are everywhere because they emerge from a fundamental tension in the human condition.'' As humans, we are a group species made up of many individuals who, in order to survive, need to organize with others. For instance, although we are born entirely dependent on others, even as we grow into independence our ability to survive (i.e., securing shelter, food, and a sense of purpose) remains deeply tied to collaboration with others \citep{Halevy-2011-Functional}. Yet this profound cooperative interdependence also gives rise to an equally fundamental competitive interdependence, as individuals must negotiate the terms of their collaboration, deciding how their relationship will be structured and how the rewards of their joint efforts will be allocated. Status hierarchies can thus be an invention to manage social situations that are characterized by cooperative interdependence to achieve valued goals and competitive interdependence to maximize individual outcomes \citep{Ridgeway-2019-Status}.

\subsection{What is Hierarchy?}

Status hierarchies are a nearly universal aspect of social organization that extend from humans to non-human animals, where fewer reside on the top than the bottom \citep{Halevy-2011-Functional, Magee-2008-SocialHierarchy, Sapolsky-2005-Influence}. Even with the wide variety of organizational forms \citep{Powell-1990-Neither, Carroll-2004-Demography} and the intentional practices and cultures designed to resist or minimize hierarchy \citep{Rothschild-1979-Collectivist, Morand-2001-Processes}, hierarchies persistently reemerge despite these efforts \citep{Leavitt-2005-Top, Tannenbaum-1974-Hierarchy}. \citet{Magee-2008-SocialHierarchy} argue that for hierarchical social relations to emerge, group members must either establish a formal system with ranked roles or participate in informal interactions through which individuals or groups become ranked along at least one valued social dimension. In this way, both formal and informal hierarchies take shape within groups and organizations.

\paragraph{Formal Hierarchy.}

As organizations grow in complexity, they increasingly formalize their hierarchies through visible markers such as job titles, reporting relationships, and organizational charts. An organizational chart, for example, provides a visual map of differentiated roles, typically depicting a small senior leadership team at the top, one or more layers of middle management, and a broad base of front-line employees responsible for daily operations and managerial support \citep{Mintzberg-1979-Structuring}. Within organizations, positions of higher formal rank are generally perceived as more valuable. Though the specific sources of this value are not always explicitly articulated, they typically encompass control over resources and deference from subordinates. Furthermore, under the assumption of effective human resource practices, higher-ranking individuals are expected to demonstrate superior skills, abilities, and motivation relative to those at lower levels, thereby conferring considerable legitimacy upon the formal hierarchy in members' eyes. Despite constant organizational flux (with individuals joining, departing, moving laterally, or ascending to higher positions), the underlying hierarchical structure itself endures beyond these individual transitions. In this respect, hierarchy exhibits stability. Once established, formal hierarchies prove relatively resistant to change, primarily because restructuring entails substantial costs.

\paragraph{Informal Hierarchy.}

Hierarchy is not solely a formal construct. Extensive research demonstrates that informal hierarchical distinctions within groups often emerge quickly and naturally \citep{Anderson-2001-Who, Bales-1951-Channels, Berger-1980-Status, Mast-2002-Dominance}. Individuals form inferences and make judgments about others' competence and power based on mere seconds of observation \citep{Ambady-1993-Half, Magee-2009-Seeing, Todorov-2005-Inferences}. Consequently, differences in task participation that emerge within minutes of interaction \citep{Fisek-1970-Process} can produce hierarchical differentiation that shapes the entire group experience. Group members also tend to exhibit high agreement about each individual's rank \citep{Mast-2004-Who}, suggesting that hierarchical differentiation is meaningful to participants even when rank ordering is based on features as subtle as nonverbal behavior.

The reasons for informal hierarchy formation in social groups vary widely. Once a particular trait or resource is deemed significant within a group or organization, members will instinctively and spontaneously establish a hierarchical ordering based on that attribute \citep{Magee-2008-SocialHierarchy}. For instance, in groups requiring minimal coordination, those who speak assertively tend to gain more status than those who speak tentatively; however, in groups requiring high coordination, the reverse pattern occurs \citep{Fragale-2006-Power}. Informal hierarchies can also arise from preexisting, stereotype-driven expectations that individuals hold about others before any direct interaction takes place \citep{Berger-1977-Status}. Race, ethnicity, gender, and class carry broad social significance, shaping interactions among group members and becoming key axes of hierarchical differentiation \citep{Berger-1977-Status, Ridgeway-1991-Social, Ridgeway-1998-HowDoStatusBeliefs}.

\subsection{Functions of Hierarchy}

The ubiquity of hierarchy suggests it fulfills important social and organizational functions. \citet{Magee-2008-SocialHierarchy} identify two primary functions: establishing social order to facilitate coordination, and providing incentives for individuals to pursue higher rank. Hierarchical structures appeal psychologically because they satisfy needs for stability and predictability \citep{Frenkel-1949-Intolerance, Sorrentino-1986-Uncertainty}, while proving organizationally effective by enabling efficient coordination. What distinguishes hierarchy from egalitarian alternatives that can also create order \citep{Krackhardt-1999-WhetherCloseOrFar} is its superior capacity for coordination through clearly defined patterns of authority and subordination. Weber's \citeyearpar{Weber-1946-Essays} analysis of bureaucracy demonstrates how hierarchy emerges as a functional adaptation to modern work demands: specialized positions are connected through vertical relationships that define behavioral expectations for leaders and followers \citep{Biggart-1984-Power, Dornbusch-1975-Authority}, enabling coordinated action. When hierarchical relationships lack clarity, work becomes inefficient and coordination suffers \citep{Greer-2007-HighPowerTeams, Overbeck-2005Internal}—even groups of high performers prove less effective without clear differentiation \citep{Groysberg-2011-TooManyCooks}.

Beyond coordination, hierarchy motivates individuals to pursue advancement by conferring material and psychological benefits at higher ranks \citep{Tannenbaum-1974-Hierarchy}, including enhanced autonomy \citep{Deci-1987-Support, Porter-1962-JobAttitudes}, internal locus of control \citep{Rotter-1966-GeneralizedExpectancies}, and power \citep{McClelland-1975-Power, Winter-1973-Power}. Weber \citeyearpar{Weber-1946-Essays} described how bureaucratic organizations create career pathways through successive promotions, and research on these internal labor markets demonstrates that advancement opportunities motivate lower-ranking members to intensify efforts toward organizational objectives \citep{Baron-1986-Structure, Pfeffer-1984-Determinants}. When promotion criteria align with organizational goals, individual ambitions serve institutional priorities, allowing hierarchy's motivational dynamics to benefit the organization. These coordination and motivation functions reveal why hierarchy has become a prevailing form of social organization: it enables groups and institutions to endure and thrive.


\subsection{Functions of Hierarchy}

The ubiquity of hierarchy suggests that it fulfills important social and organizational functions. \citet{Magee-2008-SocialHierarchy} describe two primary functions of social hierarchy. First, hierarchy establishes social order and facilitates coordination. Hierarchical structures are psychologically appealing because they satisfy individual needs for stability and are organizationally effective because they enable efficient activity coordination. Second, hierarchy provides incentives for individuals within groups and organizations. Individuals are motivated to attain higher rank to fulfill material self-interest and their desire for control. This motivation serves organizational interests insofar as rank is determined by dimensions related to group or organizational performance.

Hierarchical structures offer solutions to fundamental challenges that arise when organizing groups of individuals pursuing shared objectives. As a form of social governance, hierarchy serves as a potent remedy for uncertainty and disorder \citep{Durkheim-1997-Division, Hogg-2001-Social, Marx-1992-Early}. Through the establishment of social order, hierarchy addresses a fundamental set of human psychological needs centered on the desire for predictability, organization, and constancy \citep{Frenkel-1949-Intolerance, Sorrentino-1986-Uncertainty}. However, the psychological appeal of hierarchy's capacity to create order does not fully account for why hierarchy is favored over alternative forms of social organization. Indeed, egalitarian and balanced social arrangements can also generate order \citep{Krackhardt-1999-WhetherCloseOrFar}. What distinguishes hierarchical order is its superior effectiveness in enabling coordination among group members.

In its role as a coordination mechanism, hierarchy establishes distinct patterns of authority and subordination that enhance action coordination for numerous task types, particularly relative to more egalitarian arrangements. \citeauthor{Weber-1946-Essays}'s \citeyearpar{Weber-1946-Essays} analysis of bureaucracy indicates that hierarchy emerges as a functional adaptation to the demands of modern work. Bureaucracies distribute labor across employees \citep{Stinchcombe-1974-Industrial}, with each specialized position in this division connected through hierarchical relationships. These vertically differentiated roles define behavioral expectations for both leaders and followers \citep{Biggart-1984-Power, Dornbusch-1975-Authority}, and these role definitions enable coordinated action. When roles and hierarchical relationships lack clarity, work becomes confusing, inefficient, and frustrating, ultimately undermining coordination \citep{Greer-2007-HighPowerTeams, Overbeck-2005Internal}. Indeed, even when average task-related competence is high but clear differentiation is absent (such as when a group comprises numerous high performers), the coordination benefits of hierarchy are compromised, reducing group effectiveness and efficiency \citep{Groysberg-2011-TooManyCooks}.

Hierarchy also fulfills a motivational role by incentivizing individuals to pursue advancement within their groups and organizations, as higher positions confer greater material and psychological benefits \citep{Tannenbaum-1974-Hierarchy}. Beyond establishing order and stability, attaining elevated rank provides enhanced opportunities to satisfy control-related desires, including autonomy \citep{Deci-1987-Support, Porter-1962-JobAttitudes}, internal locus of control \citep{Rotter-1966-GeneralizedExpectancies}, and power \citep{McClelland-1975-Power, Winter-1973-Power}. This motivational function typically serves organizational interests. Weber \citeyearpar{Weber-1946-Essays} described how bureaucratic organizations create career pathways for employees by offering opportunities for successive promotions through formal hierarchical levels. Research on these career structures, often termed internal labor markets, demonstrates that the possibility of advancement motivates lower-ranking members to intensify their efforts toward organizational objectives \citep{Baron-1986-Structure, Pfeffer-1984-Determinants}. Consequently, when promotion criteria are tightly linked to organizational goals, individual ambitions align with organizational priorities, allowing the motivational dynamics of formal hierarchy to benefit the organization. Recognizing the functions hierarchy serves reveals why it has become a prevailing form of social organization: it enables groups and institutions to endure and thrive.

\subsection{Expectation States Theory}

Expectation states theory emerged from sociological research on small groups in the 1960s and 1970s, developed to explain the formation and persistence of status hierarchies in task-focused groups \citep{Berger-1974-ExpectationStatesTheory}. The theory was initially formulated to address a puzzling empirical observation: even when groups are composed of initially equal members working on collective tasks, behavioral inequalities in participation and influence rapidly emerge and stabilize. These inequalities appeared systematically linked to external status characteristics such as gender, race, and occupation, even when these characteristics had no logical relevance to the task at hand.

The theory argues that power and prestige in task-oriented interaction is determined by the performance expectations formed for one interactant compared to another \citep{Berger-1974-ExpectationStatesTheory}. A performance expectation represents an anticipation regarding the likely value of an individual's task contributions, essentially reflecting perceived task competence. When one group member is held in higher performance expectation relative to another, that individual receives more speaking opportunities, displays more assertive and confident nonverbal communication, offers more task-related suggestions, has those suggestions evaluated more favorably, and exerts greater influence \citep{Berger-1974-ExpectationStatesTheory, Ridgeway-1985-NonverbalCuesAndStatus, Ridgeway-1987-Nonverbal}. Through this process, rank-ordered performance expectations for oneself and others shape interactants' behavior in ways that validate these initial expectations. This dynamic generates a behavioral hierarchy of power and prestige that mirrors the underlying ordering of performance expectations.

Expectation states theory posits that gender influences interactional behavior through its impact on the performance expectations formed for women versus men \citep{Berger-1977-Status, Berger-1980-Status}. Gender functions as a status characteristic in society, with prevailing cultural beliefs attributing greater esteem, honor, and significance to men relative to women. Similar to other status characteristics such as race or occupation, gender status carries assumptions that individuals possessing the higher-valued state (men) demonstrate broader competence compared to those with the lower-valued state (women). When gender becomes salient in a situation, whether through compositional differences among members or through task relevance, these generalized competence assumptions become activated and shape the performance expectations formed for men and women. Once activated, gender status influences performance expectations even when unrelated to the task at hand, though its effect is amplified when the task itself is gender-typed.

Beyond status characteristics, performance expectations are shaped by interactants' established reputations for particular competencies, the compensation or rewards they receive, performance feedback, and their situational behavior \citep{Berger-1974-ExpectationStatesTheory, Berger-1985-StatusRewards}. These factors, each weighted according to task relevance, combine to form aggregate performance expectations for interactants. Consequently, the influence of gender status on power and prestige is moderated by other factors simultaneously operating within the specific context.

\subsection{Emergent Behaviors in Language Models}

Artificial Intelligence (AI) refers to advanced machine-based systems developed for application to achieve given goals or answer given questions \citep{Bengio-2024-International, Bostrom-2014-Superintelligence}. AI is a broad and quickly evolving field of study that has moved from a distant science fiction to a close reality for close to a billion people \citep{Handa-2025-WhichEconomic, Chatterji-2025-HowPeopleUse}.

The most common AI systems people interact with can be categorized as general-purpose AI. This differs from so-called `narrow AI,' a kind of AI that specialized in a tight grouping of possible tasks like playing the adversarial strategy game \textit{Go} or the recommendation algorithm that powers your social media feed \citep{Silver-2016-MasteringGo, Narayanan-2023-Understanding}. General-purpose AI rely on deep learning \citep{LeCun-2015-DeepLearning}, a method that trains artificial neural networks with many interconnected layers of nodes, enabling computers to learn from complex, unlabeled data. These networks are loosely inspired by the structure of biological networks in the brain. Most general-purpose AI rely on the transformer neural network architecture \citep{Vaswani-2017-Attention}, which has proven efficient at converting increasingly large amounts of training data and computational power into better model performance.

Development and deployment of AI systems generally follow a series of distinct steps: pre-training, fine-tuning, system integration, deployment, and post-training. Both pre-training and fine-tuning are ways of ‘training’ a general-purpose AI model. Pre-training and fine-tuning are both forms of adapting a general-purpose AI model for a specific task. In this process, the model is provided with input data and tasked with predicting related output data. For instance, it may be shown the first 500 words of a Wikipedia article and asked to generate the 501st word. At the outset, its predictions are essentially random, but as it is exposed to more data, the model adjusts based on its errors and progressively improves its accuracy.

Because AI systems, specifically language models, are trained on vast amounts of data, it is plausible that data can be unknowingly embedded with human behaviors. Developing research has begun enumerating emergent behaviors in language models including anchoring bias \citep{Lou-2024-Anchoring}, framing effects \citep{Lior-2025-WildFrame}, loss aversion \citep{Jia-2024-Decision}, escalation of commitment \citep{Barkett-2025-GettingOut}, social desirability bias \citep{Salecha-2024-Large}, truth-bias \citep{Barkett-2025-Reasoning, Markowitz-2023-Generative}, and recency bias \citep{Li-2024-NeedleBench}. 

This growing interest has sparked a closer examination into not only \textit{what} models do, but also \textit{why} they do them. The former is relatively simple to show, as experimental designs originally created for human participants can be adapted for language models \citep{Aher-2022-UsingLLMs, Anthis-2025-LLM}. However, the latter is incredibly complicated and has given rise to the field of mechanistic interpretability, which seeks to analyze language models by reverse engineering the knowledge and computational strategies they rely on, a process that more closely resembles approaches in biology or neuroscience than in traditional computer science \citep{Sharkey-2025-OpenProblemsInMechTerp, Bereska-2024-Mechanistic, McAleese-2025-MechTerp}.

\subsection{AI Alignment and Behavioral Concerns}

AI alignment refers to the challenge of ensuring that AI systems pursue objectives and exhibit behaviors that accord with human values and intentions \citep{Christian-2020-AlignmentProblem, Gabriel-2020-ArtificialIntelligenceValues}. As AI systems become increasingly capable and are deployed across more consequential domains, the importance of alignment has grown correspondingly urgent. The core difficulty lies in the fact that specifying human values with sufficient precision to guide AI behavior proves remarkably challenging, and even well-intentioned design choices during training may produce unintended behavioral patterns as systems scale in capability \citep{Amodei-2016-Concrete, Ngo-2022-AlignmentProblemDeepLearning}.

One primary concern centers on misalignment risks: the possibility that AI systems develop goals or behavioral tendencies that diverge from their intended purposes. These risks manifest in several forms. AI systems may engage in reward hacking, exploiting technical loopholes in their training objectives to achieve high rewards while failing to accomplish the intended task \citep{Pan-2022-Effects, Lehman-2018-Surprising, Schoen-2025-Stress}. Systems may also exhibit sycophantic behavior, producing outputs designed to please users rather than provide accurate or helpful information \citep{Sharma-2023-TowardsSycophancy, Perez-2022-Discovering, Malmqvist-2024-Sycophancy}. More concerning still, research has documented instances where AI systems engage in deceptive behaviors to achieve their objectives, such as providing false explanations for their actions or concealing their true capabilities \citep{Park-2023-AIDeception, Scheurer-2023-LLMDeceive, Schoen-2025-Stress, Meinke-2024-Frontier}.

Beyond deception, emerging research suggests that AI systems may exhibit power-seeking tendencies under certain conditions. Experimental studies have found that language models sometimes pursue strategies that increase their control over resources or decision-making processes, particularly when such strategies align with achieving specified goals \citep{Perez-2022-Discovering, Turner-2022-Parametrically}. These findings raise questions about whether more capable systems might develop instrumental goals related to self-preservation, resource acquisition, or influence maximization, even when these goals were not explicitly programmed \citep{Bostrom-2014-Superintelligence, Carlsmith-2022-IsPowerSeeking, Aschenbrenner-2024-SituationalAwareness, Kokotajlo-2025-ai2027, Kurzweil-2024-Singularity}.

Scalability concerns amplify these alignment challenges. Behaviors that appear benign or manageable in smaller models may become more pronounced or problematic as systems scale in size and capability \citep{Ganguli-2022-Predictability, Sharma-2023-TowardsSycophancy}. Emergent capabilities, those that appear suddenly at certain scales rather than developing gradually, make it difficult to predict which behaviors will manifest in future, more capable systems \citep{Wei-2022-Emergent}. Furthermore, some research suggests that certain learned behaviors may prove resistant to post-training interventions designed to correct them, particularly if these behaviors are deeply embedded in the model's representations \citep{Zou-2023-Universal}.

% The intersection of alignment concerns with hierarchical behavior presents particularly salient questions. If AI systems have learned hierarchical patterns from their training data, which is saturated with human status dynamics, these patterns may manifest in ways that affect human-AI collaboration and multi-agent AI coordination. An AI system that has internalized human status-seeking behaviors might, for instance, resist deferring to human judgment in collaborative settings, compete for influence in organizational contexts, or replicate human biases related to gender, race, or other status characteristics. Given that hierarchical behaviors in humans are often subtle, automatic, and resistant to explicit control, understanding whether and how AI systems exhibit analogous patterns becomes essential for anticipating and addressing potential misalignment risks.

\subsection{Literature Gap}

Despite extensive research documenting how humans navigate status hierarchies and growing evidence of emergent behavioral patterns in AI systems, a critical gap remains: we lack systematic understanding of whether and how AI language models replicate hierarchical behaviors learned from human-generated training data. This gap is particularly striking given that data used to train language models is likely saturated with hierarchical dynamics. For instance, every organizational email, news article, literary work, and social media exchange in the training corpus reflects and reinforces human status distinctions, power relations, and differential patterns of deference and assertion.

The theoretical frameworks reviewed earlier, particularly expectation states theory, provide well-established tools for understanding how humans form performance expectations, differentiate status, and enact hierarchical behaviors in group settings. Similarly, research on organizational hierarchy demonstrates how formal and informal status systems shape interaction patterns, influence distribution, and coordination. Yet these frameworks have not been systematically applied to examine AI systems. While recent work has identified specific cognitive biases and behavioral tendencies in language models, the domain of hierarchical and status-related behaviors remains unexplored.

This gap carries significant implications across multiple dimensions. From an AI safety perspective, if language models have internalized hierarchical behaviors from training data, these patterns could fundamentally shape human-AI interaction. An AI system that has learned to enact status-seeking behaviors might resist deferring to human judgment in collaborative decision-making, compete for influence in organizational contexts, or subtly reinforce harmful status hierarchies based on gender, race, or other social characteristics. Given that hierarchical behaviors in humans often operate automatically and resist conscious control \citep{Ridgeway-2019-Status}, analogous patterns in AI systems might prove similarly resistant to post-training interventions designed to promote alignment with human values.

Furthermore, as AI systems are increasingly deployed in multi-agent configurations, where multiple AI systems interact with each other and with humans, understanding hierarchical dynamics becomes essential for predicting emergent coordination patterns. Will AI agents spontaneously differentiate into hierarchies as human groups do? Will they replicate human patterns of dominance and deference, or develop novel forms of social organization? These questions bear directly on the scalability and safety of AI systems operating in complex social environments.

The theoretical sophistication of expectation states theory makes it particularly valuable for investigating these questions. The theory offers specific, falsifiable predictions about how status characteristics influence interaction patterns, how performance expectations form and propagate, and how hierarchical differentiation emerges even in initially egalitarian groups. These predictions can be operationalized through experimental paradigms that have been refined over decades of research with human participants \citep{Berger-1974-ExpectationStatesTheory, Berger-1980-Status}. Adapting these paradigms for AI systems allows for direct comparison between human and AI hierarchical dynamics, providing insights into both the extent to which AI systems have learned human social patterns and the mechanisms underlying these behaviors.

This study addresses this gap by systematically examining whether AI language models exhibit hierarchical behaviors consistent with predictions from expectation states theory. Specifically, I ask the following research question: \textit{Do language models pursue status and compete with each other for higher social standing in completing a group task? And if so, under what conditions do language models begin to compete for higher social standing?}

%\begin{enumerate}
%    \item RQ1: Do AI language models exhibit hierarchical differentiation in task-oriented interactions, such that they demonstrate patterns of dominance, deference, and influence that mirror those predicted by expectation states theory for human groups?
%    \item RQ2: Do status characteristics such as gender and occupation influence AI systems' hierarchical behaviors in ways consistent with their effects on human interaction, even when these characteristics have no logical relevance to task performance?
%    \item RQ3: Are AI systems' hierarchical behaviors stable across contexts and resistant to explicit instructions to behave egalitarianly, similar to the persistence of status effects in human groups?
%\end{enumerate}

By answering these questions, this research makes several contributions. First, it extends our understanding of emergent social behaviors in AI systems beyond individual cognitive biases to complex, socially embedded patterns of interaction. Second, it demonstrates the value of applying established social psychological theory to AI behavior, opening new avenues for predicting and understanding how these systems will function in social contexts. Third, it provides empirical evidence relevant to AI alignment concerns, particularly regarding whether hierarchical behaviors learned from training data persist despite fine-tuning efforts designed to make systems helpful, harmless, and honest. Finally, it establishes a methodological foundation for future research examining other dimensions of social behavior in AI systems, from coalition formation to norm enforcement to collective decision-making.